% Part: first-order-logic
% Chapter: axiomatic-deduction
% Section: proving-things

\documentclass[../../../include/open-logic-section]{subfiles}

\begin{document}

\iftag{FOL}
      {\olfileid{fol}{axd}{pro}}
      {\olfileid{pl}{axd}{pro}}

\olsection{\usetoken{P}{derivation}示例}

\begin{ex}
假设我们要证明 $(\lnot !D \lor !E) \lif (!D \lif
!E)$。显然，它不是我们任何公理的实例，因此我们必须使用 \MP{} 规则来推导它。我们唯一的规则是~\MP，在给定 $!A$ 和 $!A \lif !B$ 的情况下，它允许我们得出~$!B$。一种策略是使用 \olref[prp]{ax:lor3}，令 $!A$ 为 $\lnot !D$、$!B$ 为 $!E$、$!C$ 为 $!D \lif !E$，即实例
\[
(\lnot !D \lif (!D \lif !E)) \lif ((!E \lif (!D \lif !E)) \lif ((\lnot
!D \lor !E) \lif (!D \lif !E))).
\]
为什么？两次应用 MP 会得到最后那部分，而这正是我们想要的。我们很容易看出 $\lnot !D \lif (!D \lif !E)$ 是 \olref[prp]{ax:lnot2} 的一个实例，而 $!E \lif (!D \lif !E)$ 是 \olref[prp]{ax:lif1} 的一个实例。因此我们的推导是：
\begin{derivation}
  1. &  $\lnot !D \lif (!D \lif !E)$ & \olref[prp]{ax:lnot2} \\
  2. & $(\lnot !D \lif (!D \lif !E)) \lif {}$\\
  &\qquad $((!E \lif (!D \lif !E)) \lif ((\lnot !D \lor !E) \lif (!D \lif !E)))$ & \olref[prp]{ax:lor3}\\
  3. & $(!E \lif (!D \lif !E)) \lif ((\lnot !D \lor !E) \lif (!D \lif !E))$ & 1, 2, \MP\\
  4. & $!E \lif (!D \lif !E)$ & \olref[prp]{ax:lif1}\\
  5. & $(\lnot !D \lor !E) \lif (!D \lif !E)$ & 3, 4, \MP
\end{derivation}
\end{ex}

\begin{ex}\ollabel{ex:identity}
让我们尝试寻找 $!D \lif !D$ 的一个推导。它不是公理的实例，所以我们必须用 \MP 来推导它。
\olref[prp]{ax:lif1} 是形如~$!A \lif !B$ 的公理，我们可以对其应用~\MP。当然，为了使推导有意义，在这种情况下 \MP 作为正确步骤所验证的 $!B$ 必须是~$!D \lif !D$，因为这是我们想要推导的。这意味着 $!A$ 也必须是 $!D$，也就是说，我们可以考虑 \olref[prp]{ax:lif1} 的如下实例：
\[
!D \lif (!D \lif !D)
\]
为了应用 \MP，我们还需要证明对应的第二个前提，即~$!A$。但在我们的情形中，那就是~$!D$，而我们无法单独推导出~$!D$。所以我们需要一个不同的策略。

另一个只涉及~$\lif$ 的公理是 \olref[prp]{ax:lif2}，即
\[
(!A \lif (!B \lif !C)) \lif ((!A \lif !B) \lif (!A \lif !C))
\]
我们可以通过两次应用 \MP 来得到最后那个嵌套的条件句。同样，这意味着我们希望 \olref[prp]{ax:lif2} 有一个实例，其中 $!A \lif !C$ 是 $!D \lif !D$，即我们的目标公式。那么当然，$!A$ 和 $!C$ 都是~$!D$。我们应该如何选择~$!B$，使得 $!A \lif (!B \lif !C)$ 和 $!A \lif !B$（即在我们的例子中的 $!D \lif (!B \lif !D)$ 和 $!D \lif !B$）也都是可推导的呢？其实，无论 $!B$ 取什么，第一个公式都已经是 \olref[prp]{ax:lif1} 的一个实例。而如果 $!B$ 是 $(!D \lif !D)$，那么 $!D \lif !B$ 就是 \olref[prp]{ax:lif1} 的另一个实例。因此，我们的推导是：
\begin{derivation}
  1. & $!D \lif ((!D \lif !D) \lif !D)$ & \olref[prp]{ax:lif1}\\
  2. & $(!D \lif ((!D \lif !D) \lif !D)) \lif {}$\\
  & \qquad $((!D \lif (!D \lif !D)) \lif (!D \lif !D))$ & \olref[prp]{ax:lif2}\\
  3. & $(!D \lif (!D \lif !D)) \lif (!D \lif !D)$ & 1, 2, \MP\\
  4. & $!D \lif (!D \lif !D)$ & \olref[prp]{ax:lif1}\\
  5. & $!D \lif !D$ & 3, 4, \MP
\end{derivation}
\end{ex}

\begin{ex}\ollabel{ex:chain}
有时我们想要证明，从其他一些公式~$\Gamma$ 可以推导出某个公式。例如，我们来证明从 $\Gamma = \{!A \lif !B,
!B \lif !C\}$ 可以推导出 $!A \lif !C$。
\begin{derivation}
  1. & $!A \lif !B$ & \Hyp\\
  2. & $!B \lif !C$ & \Hyp\\
  3. & $(!B \lif !C) \lif (!A \lif (!B \lif !C))$ & \olref[prp]{ax:lif1} \\
  4. & $!A \lif (!B \lif !C)$ & 2, 3, \MP\\
  5. & $(!A \lif (!B \lif !C)) \lif {}$\\
  & \qquad $((!A \lif !B) \lif (!A \lif !C))$ & \olref[prp]{ax:lif2}\\
  6. & $((!A \lif !B) \lif (!A \lif !C))$ & 4, 5, \MP\\
  7. & $!A \lif !C$ & 1, 6, \MP
\end{derivation}
标记为“\Hyp”（表示“假设”）的行，表示该行上的公式属于~$\Gamma$。
\end{ex}

\begin{prop}
  \ollabel{prop:chain} 如果 $\Gamma \Proves !A \lif !B$ 且 $\Gamma
  \Proves !B \lif !C$，那么 $\Gamma \Proves !A \lif !C$。
\end{prop}

\begin{proof}
  假设 $\Gamma \Proves !A \lif !B$ 且 $\Gamma \Proves !B \lif
  !C$。则存在从~$\Gamma$ 到 $!A \lif !B$ 的一个推导，以及从~$\Gamma$ 到 $!B \lif !C$ 的一个推导。将这两个推导拼接起来，合并为一个推导。现在添加上例中推导的第 3--7 行。这就得到了一个从~$\Gamma$ 到 $!A \lif !C$ 的推导——$!A \lif !C$ 也是新推导的最后一行。注意，如果将对第 1 行的引用替换为对 $!A \lif !B$ 推导最后一行的引用，将对第 2 行的引用替换为对 $!B \lif !C$ 推导最后一行的引用，那么第 4 行和第 7 行的依据仍然有效。
\end{proof}

\begin{prob} 
通过展示从公理出发的推导，证明以下命题：
\begin{enumerate} 
\item $(!A \land !B) \lif (!B \land !A)$
\item $((!A \land !B) \lif !C) \lif (!A \lif (!B \lif !C))$
\item $\lnot(!A \lor !B) \lif \lnot !A$
\end{enumerate} 
\end{prob}

\end{document}